\begin{textsample}{POS Dim 1 – 2011\_c4 – Score -2.00 – 2011\_c4/enade\_2011\_art\_vis\_14aea27cb0.txt}  \label{ex:f1_pos_321}
CAPÍTULO 4 PERCEPÇÃO DA PROVA As análises feitas neste capítulo tratam das percepções dos concluintes da Área de Artes Visuais ( Licenciatura ) sobre a prova aplicada no ENADE/2011.
Estas percepções foram mensuradas por meio de nove questões que avaliaram desde o grau de dificuldade da prova até o tempo gasto para concluí -la.
As percepções sobre a prova foram relacionadas com o desempenho dos estudantes e com a Grande Região de funcionamento do curso.
O questionário de percepção da prova encontra -se ao final do Anexo IV, que traz a reprodução da prova.
O desempenho dos estudantes foi classificado em quatro quartos.
Para tanto, esse desempenho foi ordenado de forma ascendente.
O percentil 25, P25, também conhecido como primeiro quartil, é a nota de desempenho que deixa um quarto ( 25 \% ) dos valores observados abaixo e três quartos acima.
A Figura 1 apresenta uma ilustração deste conceito.
O quarto inferior de desempenho é composto pelas notas abaixo do primeiro quartil.
Já o percentil 75, P75, também conhecido como terceiro quartil, é o valor para o qual há três quartos ( 75 \% ) dos dados abaixo e um quarto acima dele.
O quarto superior de desempenho é composto pelas notas iguais ou acima do terceiro quartil.
O percentil 50, P50, também conhecido como mediana, é o valor que divide as notas em dois conjuntos de igual tamanho.
O segundo quarto inclui valores entre o primeiro quartil ( P25 ) e a mediana.
O terceiro quarto contém os valores entre a mediana ( P50 ) e o terceiro quartil ( P75 ).
Vale ressaltar que percentis, quartis e medianas são pontos que não obrigatoriamente pertencem ao conjunto original de dados, ao passo que os quartos são subconjuntos dos dados originais.
A seguir, serão apresentados gráficos com resultados selecionados, relativos às nove questões avaliadas por grupos de estudantes.
Os gráficos apresentam nas barras o percentual de alunos que assinalaram uma das opções ou a soma das porcentagens daqueles que assinalaram duas ( ou três ) delas.
Por exemplo, para as questões 1 e 2, os gráficos apresentam a porcentagem total de participantes que assinalaram as opções ( D ) difícil e ( E ) muito difícil.
Em cada barra foram assinalados também os extremos do intervalo de confiança de 95 \% como linhas verticais unidas por uma linha horizontal na forma da letra H maiúscula, semelhantemente aos gráficos do Capítulo 3.
As Tabelas no Anexo II apresentam os valores absolutos e a distribuição percentual das alternativas válidas das nove questões, segundo o mesmo recorte de desempenho dos alunos e Grande Região de funcionamento do curso. 4.1 GRAU DE DIFICULDADE DA PROVA 4.1.1 Componente de Formação Geral Ao avaliarem “ Qual o grau de dificuldade desta prova na parte de Formação Geral? ” ( Questão 1 ), 29,5 \% do grupo de inscritos e presentes optaram pelas alternativas difícil ou muito difícil.
Entretanto, para mais da metade dos estudantes ( 59,8 \% ), o Componente de Formação Geral da prova foi considerado com grau de dificuldade médio ( Gráfico 4.1, Gráfico 4.2 e, no Anexo II, a Tabela II.1 ).
O percentual de estudantes que consideraram a prova como difícil ou muito difícil foi maior na região Centro-Oeste, onde a proporção foi de 34,7 \%, enquanto a de menor incidência foi a Nordeste, com 11,2 \%.
No Gráfico 4.1 é possível observar que esta diferença é estatisticamente significativa.
Nas Grandes Regiões, a proporção de presentes à prova que consideraram o Componente de Formação Geral como sendo de grau de dificuldade médio esteve entre 59,0 \% na região Sudeste e 72,3 \% na região Norte.
O percentual de alunos que consideraram a prova difícil ou muito difícil foi decrescente segundo desempenho, com diferenças estatisticamente significativas.
Em cada um dos quartos de desempenho a proporção de alunos que consideraram a prova difícil ou muito difícil foi : 44,8 \% no primeiro quarto, 35,2 no segundo, 25,3 no terceiro e 14,0 \% no último quarto.
Para todos os quartos de desempenho a alternativa modal para esta pergunta foi médio, variando de 45,7 \% no primeiro quarto a 68,5 \% dos respondentes no quarto que \textbf{contam} as maiores notas. 4.1.2 Componente de Conhecimento Específico Ao responderem à Questão 2 – “ Qual o grau de dificuldade desta prova na parte de Componente Específico? ” – 31,6 \% do grupo de estudantes classificaram -na como difícil ou muito difícil.
Além disso, o Componente de Conhecimento Específico da prova foi considerado com grau de dificuldade médio por 59,0 \% dos alunos ( Gráfico 4.3, Gráfico 4.4, e, no Anexo II, a Tabela II.2 ).
A análise das respostas dos estudantes quanto ao grau de dificuldade do Componente de Conhecimento Específico da prova, agregado por Grande Região, mostra que a diferença entre a maior e a menor proporção de alunos que a avaliaram como difícil ou muito difícil é estatisticamente significativa : 21,1 \% na região Nordeste e 37,6 \% na Centro-Oeste.
O percentual de alunos que classificaram o grau de dificuldade como médio, no Componente de Conhecimento Específico, variou de 54,4 \% a 67,7 \%, para as mesmas regiões, Centro-Oeste e Nordeste, respectivamente.
Considerando -se a avaliação da dificuldade das questões do Componente de Conhecimento Específico da prova, de acordo com o desempenho dos estudantes, se observa diferença estatisticamente significativa entre os quartos extremos.
A proporção dos que classificaram a parte específica como difícil ou muito difícil variou de 24,0 \% ( 4º quarto ) a 37,4 \% ( 1º quarto ).
Já a alternativa modal para a Questão 2 foi o grau médio, com 51,0 \% do quarto inferior e 65,4 \% do superior optando por esta resposta. 4.2 EXTENSÃO DA PROVA EM RELAÇÃO AO TEMPO TOTAL Indagados quanto à extensão da prova, em relação ao tempo total oferecido para a sua resolução ( Questão 3 ), os estudantes apontaram, com maior incidência, a alternativa que considerava a extensão adequada, para todas as agregações consideradas.
A única exceção ocorreu na região Nordeste, onde a escolha da opção muito longa feita por 34,6 \% dos alunos superou um pouco os 32,3 \% que escolheram a opção adequada ( Gráfico 4.5, Gráfico 4.6, e, no Anexo II, a Tabela II.3 ).
No Brasil como um todo, o percentual de alunos que responderam ser a extensão da prova adequada foi de 44,4 \%.
Já 51,3 \% dos inscritos presentes consideraram que o exame foi longo ou muito longo e 4,3 \% o avaliaram como curto ou muito curto.
Entre as Grandes Regiões a proporção daqueles que avaliaram a prova como longa ou muito longa em relação ao tempo total destinado à sua resolução variou de 45,0 \% na região Norte até 61,7 \% na região Nordeste.
Apesar de esta diferença parecer grande, ela não é estatisticamente significativa.
Entre as demais regiões a variação é visivelmente pequena.
Considerando -se o desempenho dos alunos, nota -se que a escolha modal para todos os quartos de desempenho foi considerar a extensão da prova adequada quanto ao tempo.
Esta opção foi escolhida por 45,2 \% e 45,4 \% dos estudantes com desempenho nos dois quartos superiores, respectivamente, terceiro e quarto.
Nos quartos inferiores a extensão da prova foi considerada adequada por 43,2 \% no primeiro quarto e 43,6 \% no segundo.
No Gráfico 4.6, pode -se constatar que para os quatro quartos de desempenho a proporção de estudantes que consideraram a prova longa ou muito longa, em relação ao tempo total destinado à sua resolução, manteve -se entre 49,9 \% e 52,4 \%, sem diferenças estatisticamente significativas. 4.3 COMPREENSÃO DOS ENUNCIADOS DAS QUESTÕES 4.3.1 Componente de Formação Geral Com relação aos enunciados das questões do Componente de Formação Geral ( Questão 4 ), as opiniões foram positivas, já que 74,7 \% dos alunos avaliados consideraram os enunciados de todas ou da maioria das questões claros e objetivos ( Gráfico 4.7, Gráfico 4.8, e, no Anexo II, a Tabela II.4 ).
Na análise regional, a percentagem de estudantes que avaliaram que todos ou a maioria dos enunciados das questões do Componente de Formação Geral estavam claros e objetivos variou de 72,6 \% na região Sudeste a 80,3 \% na região Nordeste, sendo que nenhuma das diferenças entre as Grandes Regiões é estatisticamente significativa.
A análise das percepções dos estudantes sobre a clareza e objetividade dos enunciados permite afirmar que todos ou a maioria dos enunciados de questões relativas ao Componente de Formação Geral foram considerados claros e objetivos para a maior parte dos respondentes ( maior ou igual a 72,5 \% em todas as regiões e maior ou igual a 61,9 \% para todos os quartos de desempenho ).
Segundo o desempenho, observa -se que a proporção dos que emitiram esta opinião cresce conforme o desempenho aumenta, com diferenças estatisticamente significativas.
No quarto superior, a clareza e objetividade de todos ou da maioria dos enunciados das questões foi percebida por 87,1 \% dos alunos e no quarto de desempenho inferior tal avaliação foi emitida por 61,9 \% deles. 4.3.2 Componente de Conhecimento Específico Com relação aos enunciados das questões do Componente de Conhecimento Específico da prova, para 75,9 \% dos estudantes avaliados da Área de Artes Visuais ( Licenciatura ) a clareza e a objetividade ( Questão 5 ) estavam presentes em todas ou na maioria das questões ( Gráfico 4.9, Gráfico 4.10, e no Anexo II, a Tabela II.5 ).
A maioria dos estudantes de todas as Grandes Regiões brasileiras considerou claros e objetivos todos ou a maioria dos enunciados das questões do Componente de Conhecimento Específico da prova, percentual sempre maior do que 73 \%.
Observa -se no Gráfico 4.9 que as diferenças entre Grandes Regiões não são estatisticamente significativas.
A proporção de estudantes que consideraram os enunciados das questões claros e objetivos apresenta uma tendência crescente em relação ao aumento de desempenho : mais elevada no quarto superior ( 83,2 \% ) se comparada ao quarto inferior de desempenho ( 64,3 \% ).
Esta diferença é estatisticamente significativa, bem como a diferença entre a frequência desta escolha no segundo quarto em relação aos dois quartos seguintes. 4.4 SUFICIÊNCIA DAS INFORMAÇÕES/INSTRUÇÕES FORNECIDAS Ao avaliarem as informações/instruções fornecidas para a resolução das questões ( Questão 6 ), 83,6 \% dos respondentes da Área de Artes Visuais ( Licenciatura ) de todo o Brasil afirmaram que estas eram até excessivas ou suficientes em todas ou na maioria das questões ( Gráfico 4.11, Gráfico 4.12, e, no Anexo II, a Tabela II.6 ).
Quanto à distribuição de respondentes pelas Grandes Regiões observa -se que a proporção de estudantes que consideraram as informações/instruções fornecidas até excessivas ou suficientes em todas ou na maioria das questões foi sempre superior a 80 \%, chegando a 85,9 \% na região Norte.
As diferenças entre as regiões não chegaram a ser estatisticamente significativas.
Levando -se em conta o desempenho dos participantes, notam -se diferenças estatisticamente significativas entre as opiniões de estudantes do quarto inferior e superior de desempenho, como mostra o Gráfico 4.12.
O percentual de participantes que avaliaram as informações/instruções como até excessivas ou suficientes em todas ou na maioria das questões foi mais elevado no quarto superior ( 92,3 \% ), percentual superior à média nacional ( 83,6 \% ).
Já no quarto inferior, a suficiência das informações/instruções declarada como até excessiva em todas ou na maioria das questões foi percebida por 72,3 \% dos respondentes. 4.5 DIFICULDADE ENCONTRADA AO RESPONDER À PROVA Perguntados sobre as dificuldades com as quais se depararam ao responder à prova ( Questão 7 ), 11,2 \% dos estudantes apontaram o desconhecimento do conteúdo.
Para 47,0 \%, a forma diferente de abordagem do conteúdo foi indicada como dificuldade.
Já a falta de motivação para fazer a prova foi a dificuldade apontada por 20,2 \% dos respondentes.
Considerando -se todo o Brasil, 17,4 \% dos respondentes afirmaram que não tiveram qualquer tipo de dificuldade para responder à prova ( Tabela II.7 no Anexo II ).
Os Gráficos 4.13 e 4.14 apresentam os percentuais de estudantes que apontaram o desconhecimento do conteúdo como dificuldade percebida ao responder à prova.
Na análise por Grandes Regiões, o percentual de inscritos e presentes que apontaram o desconhecimento do conteúdo como dificuldade ao responder à prova chegou a 13 \%.
Os percentuais variaram de 9,3 \% na região Norte a 13,0 \% na Nordeste, sendo esta diferença não significativa estatisticamente.
A forma diferente de abordagem do conteúdo foi a escolha modal dos estudantes, com percentuais que variaram de 39,7 \% ( região Nordeste ) a 52,4 \% ( Sul ).
O percentual de alunos que citaram a falta de motivação como dificuldade variou de 15,3 \% ( região Nordeste ) a 24,3 \% ( região Sudeste ).
Os que declararam não ter qualquer dificuldade para responder à prova variaram de 16,1 \% na região Sudeste a 22,1 \% na Nordeste.
Com relação aos quartos de desempenho, o desconhecimento do conteúdo, foi a opção escolhida por 12,2 \% dos estudantes do quarto superior e 14,4 \% do quarto inferior.
A alternativa modal para os alunos, quando agregados pelos quartos de desempenho, foi que a dificuldade encontrada foi causada pela forma diferente de abordagem do conteúdo : 48,5 \% no quarto inferior e 38,9 \% do quarto superior assim o responderam. 4.6 CONTEÚDOS DAS QUESTÕES OBJETIVAS DA PROVA Ao analisarem os conteúdos das questões objetivas da prova ( Questão 8 ), um percentual pequeno dos estudantes avaliados, apenas 8,1 \%, afirmou que não estudou ainda a maioria desses conteúdos ( Gráficos 4.15, Gráfico 4.16, e a Tabelas II.8 no Anexo II ).
A maioria ( 68,0 \% ) afirmou ter estudado e aprendido muitos ou todos os conteúdos.
Na análise por Grande Região, a proporção de respondentes que escolheram a opção não estudou ainda a maioria desses conteúdos, foi pequena.
Observa -se que nas regiões Sul ( 11,3 \% ) e Norte ( 8,2 \% ), apesar de pequenas, as proporções foram maiores do que a média nacional ( 8,1 \% ).
Não se observa diferença estatisticamente significativa entre as regiões.
Em todas as regiões, a maioria dos presentes afirmou ter estudado e aprendido muitos ou todos os conteúdos, com proporções variando entre 64,1 \% na região Sul e 73,5 \% na Nordeste.
Considerando -se separadamente as opiniões de estudantes dos quatro quartos de desempenho, observa -se que, no quarto inferior, 15,6 \% ofereceram como resposta que não estudou ainda a maioria desses conteúdos, sendo 3,3 \% os do quarto superior com a mesma resposta.
A diferença entre os alunos que optaram por este motivo de dificuldade nos quartos extremos é estatisticamente significativa.
Tendo em conta o quarto superior, 82,8 \% dos alunos afirmaram ter estudado e aprendido muitos ou todos os conteúdos. 4.7 TEMPO GASTO PARA CONCLUIR A PROVA Ao responderem sobre o tempo de conclusão da prova ( Questão 9 ), mais da metade dos estudantes ( 72,5 \% ) afirmou ter gasto entre duas e quatro horas ( Gráfico 4.17, Gráfico 4.18 e, no Anexo II, a Tabela II.9 ).
Considerando -se as cinco Grandes Regiões brasileiras, os que utilizaram entre duas e quatro horas para finalizar a prova superaram o percentual nacional nas regiões Norte ( 74,5 \% ), Sul ( 73,2 \% ) e Centro-Oeste ( 75,7 \% ).
Em quatro das cinco Grandes Regiões o percentual de alunos que dispensaram entre duas e quatro horas para concluir a prova, ficou acima de 70 \%, como mostra o Gráfico 4.17.
Uma vez sendo analisadas as alternativas escolhidas pelos estudantes que se situam nos diferentes quartos de desempenho, observa -se que uma maior proporção de participantes no quarto superior declarou ter gasto entre duas e quatro horas para concluir a prova quando comparados com os do quarto inferior, respectivamente 79,7 \% e 61,3 \%.
As diferenças entre o primeiro e quarto quartos são estatisticamente significativas e evidenciam uma tendência crescente.

% matched lemmas: contar
\end{textsample}
